Forecasting models for operational demand planning are often evaluated by numerical prediction accuracy, yet operational systems do not execute forecasts directly. They execute planning signals such as inventory targets, replenishment recommendations, staffing plans, or capacity allocations. A forecast that improves error metrics can still produce volatile or infeasible plans when downstream execution capacity, model-switching burden, inventory exposure, and planning stability are ignored. This paper formulates forecast-driven demand planning as a feasibility-constrained model selection problem. The proposed evaluation layer converts candidate forecasts into executable planning signals and compares deployable strategies using both raw operational metrics and a normalized planning utility that includes inventory cost, planning volatility, switching burden, and execution violations. Empirically, this study uses DataCo supply-chain execution data to anchor execution-risk scenarios, Favorita retail demand data as the main forecast-to-plan proof of concept, M5 retail demand data as a hierarchical robustness setting, and Walmart data as a weekly business-context robustness check. Across the current paper-facing evaluations, the accuracy-first deployable method is not the operational-loss-optimal deployable strategy in 91.1\% of 123 evaluation groups. The current results suggest that execution-risk-sensitive objectives can reorder deployable strategy rankings, without implying that any single strategy universally dominates. The implication is that operational feasibility should be treated as a first-class planning objective alongside forecast accuracy.
